% One-page explainer: intonation & vibrato via a spectral-mixture GP.
% Shareable: self-contained, no project-internal terminology.
% Compile: tectonic sm_estimator_onepager.tex
\documentclass[10pt,a4paper]{article}
\usepackage[margin=1.9cm]{geometry}
\usepackage{amsmath,amssymb}
\usepackage[T1]{fontenc}
\usepackage{microtype}
\setlength{\parindent}{0pt}
\setlength{\parskip}{4.5pt}
\pagestyle{empty}

\begin{document}

\begin{center}
{\large\bfseries Estimating intonation and vibrato with a spectral-mixture Gaussian process}\\[2pt]
{\small Raynaldi Lalang \,--\, one-page summary, September 2026}
\end{center}

\textbf{The task.}
From a recording of one played note we want three numbers, each with an honest
uncertainty: the \emph{intonation} $c$ (how sharp or flat the note's held centre
sits, in cents), the \emph{vibrato rate} $f$ (oscillations per second), and the
\emph{vibrato extent} $\gamma$ (how far the pitch swings, in cents). A pitch
tracker turns the note's audio into a \emph{cents curve}: the pitch at each
analysis frame, written as the deviation in cents from the note's written pitch.
The curve's constant part is the intonation; its oscillation is the vibrato.

\textbf{The current method and its limits.}
The standard approach fits one rigid shape to the curve by least squares: a
constant centre plus a phase-locked sinusoid. Two measured problems follow.
First, real notes \emph{drift}: the centre moves during the note (we measured a
median drift of $\sim$10 cents), and a rigid shape pushes that drift into the
residual, biasing the fit. Second, where the shape assumption breaks (short
notes, weak or irregular vibrato) the fit is unreliable, so hard rules must
discard those notes entirely --- in our corpus the majority of notes.

\textbf{The idea: state what we know as a distribution over curves.}
What we actually know about a cents curve is knowledge about \emph{frequencies}:
oscillation near the vibrato rate (5--8\,Hz) and slow variation near zero
frequency (the drift). A Gaussian process (GP) prior can express exactly this,
because choosing a GP kernel is equivalent to choosing how the curve's energy is
spread over frequencies. Wilson \& Adams (ICML 2013) model that spread directly
as a mixture of a few spectral bumps; the corresponding kernel is closed-form:
\begin{equation*}
k(\tau) \;=\; \textstyle\sum_{q} w_q\, e^{-2\pi^2 \tau^2 v_q} \cos(2\pi \tau \mu_q),
\qquad
\text{one damped cosine per bump.}
\end{equation*}
Each component $q$ has a frequency $\mu_q$, a power $w_q$, and a phase-coherence
decay $v_q$. Two components suffice here --- one near the vibrato rate, one at
zero frequency for the drift --- so the model for the observed frames
$y_j = c(t_j)$ is
\begin{equation*}
c(t) = c + g(t), \quad g \sim \mathcal{GP}(0, k), \qquad
\operatorname{Cov}(y_j, y_{j'}) =
\underbrace{w_1 e^{-2\pi^2\tau^2 v_1}\cos(2\pi\tau\mu_1)}_{\text{vibrato}}
+ \underbrace{w_2\, e^{-2\pi^2\tau^2 v_2}}_{\text{drift }(\mu_2=0)}
+ \underbrace{\sigma^2 \mathbf{1}[j{=}j']}_{\text{tracker noise}}.
\end{equation*}
The point of this particular kernel: \emph{its parameters are the quantities we
want to measure}, so fitting the model is the estimation itself. The sine fit is
its degenerate limit (one component, perfectly phase-locked, no drift term).

\textbf{How each number is estimated.}
All hyperparameters are fitted at once by maximising the exact GP marginal
likelihood of the note's frames,
$\log \mathcal{N}\!\bigl(y;\, c\mathbf{1},\, K(\theta)\bigr)$.
\emph{Intonation} $c$ is the correlation-aware average: given the kernel, the
constant mean has an exact closed-form Gaussian posterior, and because the model
knows which wiggle is vibrato and which is drift, an incomplete vibrato cycle or
a drifting tail does not bias the centre (the naive curve mean has exactly that
bias). \emph{Rate} $\mu_1$ is found where the frames' correlation pattern peaks
the likelihood: the kernel predicts that frames one vibrato period apart agree
and frames half a period apart disagree --- an oscillating correlation pattern
readable from the data \emph{without knowing the phase}. A coarse grid over
3--10\,Hz initialises the search (the likelihood is multimodal in frequency).
\emph{Extent} follows from the power: too little $w_1$ forces the oscillation to
be explained as white noise (which predicts no inter-frame correlation), too
much predicts bigger swings than observed; the optimum is the variance the
oscillation actually carries, and a sinusoid of amplitude $\gamma$ has variance
$\gamma^2/2$, so $\hat\gamma = \sqrt{2 w_1}$.

\textbf{Uncertainty, and graceful failure.}
The centre's variance is exact. Rate and extent inherit theirs from the
curvature of the log-likelihood at its optimum: many coherent cycles give a
needle-sharp peak in $\mu_1$ and a small variance; a short or irregular note
gives a nearly flat surface and a wide one. This replaces the discard rules: a
weak note is not refused, it \emph{reports itself as uninformative}, and the
uncertainty propagates into the across-note model that consumes these per-note
estimates. That model, and the set of per-note quantities it works with, are
unchanged --- this is a swap of the estimator behind three of them. (One known
limit: a stationary kernel averages phase away, so the vibrato \emph{onset
delay} cannot come from this model and keeps its current estimator.)

\textbf{How it will be judged.}
On development data only ($\sim$5{,}000 vibrato-identifiable notes from a
chamber-music corpus with ground-truth pitch), four measures. (1)~\emph{Parameter
accuracy}, the pass/fail criterion: each estimator runs on the tracked curve and
on the ground-truth curve of the same note and is scored against its own
ground-truth counterpart, so neither is judged by the other's definition of the
targets. (2)~\emph{Calibration}: do the reported uncertainties cover the truth
at their nominal rate (coverage, log score, PIT)? (3)~\emph{The refused notes}:
on notes the sine fit discards, does the GP's wide posterior cover the
ground-truth value? (4)~\emph{The curve itself}: both models' posteriors scored
directly against the ground-truth curve's frames, which requires no target
definition at all. If the GP does not beat the sine fit on (1), the idea is
dropped; any adoption beyond a development study would be preregistered and
confirmed on fresh data, and the already-confirmed results stand unchanged.

\end{document}
